\scanpage{71}
- \mspage{29} -
Since $\operatorname{Coim}f=\operatorname{Im}f$, it follows that $\operatorname{Coim}'f=\operatorname{Im}'f$.

ii) Let $0\longrightarrow D'\longrightarrow D\longrightarrow D''\longrightarrow0$ be an exact sequence of directed systems in $\underline A'$ with $D'(i)'=A_i'$, etc.
(i.e., for for each $i\in\underline D$, $0\longrightarrow A_i'\longrightarrow A_i\longrightarrow A_i''\longrightarrow$ is exact in $\underline A'$.

Then, by 4.7,
\[
*)\quad\operatorname{dir\,lim}A_i'\longrightarrow\operatorname{dir\,lim}A_i
\longrightarrow\operatorname{dir\,lim}A_i''\longrightarrow0
\]
is exact in $\underline A'$. We must show it is also exact on the left. Now, since $I_{\underline A\underline A'}$ has a left adjoint,
\[
0\longrightarrow I_{\underline A\underline A'}(A_i')
\longrightarrow I_{\underline A\underline A'}(A_i)
\longrightarrow I_{\underline A\underline A'}(A_i'')
\]
is exact in $\underline A$. (In general it is not exact on the right in $\underline A$).
Hence, since $\underline A$ is right perfect,
\[
0\longrightarrow\operatorname{dir\,lim}I_{\underline A\underline A'}(A_i')
\longrightarrow\operatorname{dir\,lim}I_{\underline A\underline A'}(A_i)
\longrightarrow\operatorname{dir\,lim}I_{\underline A\underline A'}(A_i'')
\]
is exact. Since $R$ is an exact functor, the sequence remains exact after applying $R$. But, by 4.14,
\[
R[\operatorname{dir\,lim}I_{\underline A\underline A'}(A_i')]
=\operatorname{dir\,lim}A_i'
\]
etc. Hence, the sequence *) is exact on the left.

\uline{6.3.}\ \uline{Theorem}: Let $\underline C$ be a full subcategory of a functor category $\underline F(\underline B,\underline A)$ where $\underline A$ is right complete. Suppose $R:\underline F(\underline B,\underline A)\longrightarrow\underline C$ is a right reflection. Then

i) $\underline C$ is right complete

ii) If $\underline A$ has a generator, then so does $\underline C$.

Suppose further that $\underline A$ is abelian, $\underline C$ is left closed, and $R$ is exact. Then

iii) $\underline C$ is abelian.

iv) If $\underline A$ is right perfect, then so is $\underline C$.
\mspage{70}
